\section{Estratégia Econométrica}
\label{sec:estrategia}

A comparação descritiva mostrou que o crescimento da produção de grãos coexistiu com taxas de conversão menores em 2024 que em 2001 e com o predomínio da pastagem entre as áreas abertas na amostra fixa. Essas trajetórias agregadas não revelam se municípios com maior potencial agrícola apresentam associação distinta entre preços e conversão, que constitui o contraste de H1. Para examiná-lo, a estratégia combina a variação anual das cotações com diferenças de capacidade entre municípios, considerando características permanentes e mudanças comuns aos grupos de bioma predominante. As extensões situam essa associação diante da ocupação agrícola anterior, do acesso aos mercados e dos registros administrativos, enquanto os exercícios de robustez delimitam sua estabilidade e seu alcance interpretativo.

\subsection{Especificação de referência, exposição potencial e pacote tecnológico}
\label{estrategia-especificacao}

A análise utiliza o painel anual dos 141 municípios de Mato Grosso na malha do IBGE de 2023 entre 2001 e 2024, que perfaz 3.384 pares município--ano potenciais. Como a cobertura dos mapas de transição do MapBiomas Coleção 10.1 não é integral em todos os anos, a amostra restringe-se aos pares em que a conversão da vegetação nativa é efetivamente observada \cite{MapBiomas2026}. A exclusão das 612 observações sem cobertura completa preserva 2.772 pares válidos em 137 municípios, todos com informação de precipitação. Esse conjunto elegível difere das amostras efetivas de estimação: a remoção de observações em grupos unitários de efeitos fixos resulta em 2.769 observações na linha de base e 2.767 na especificação de referência, que emprega efeitos de bioma por ano.

Para cada cultura $c \in \{s, m\}$, correspondente à soja e ao milho, a especificação de referência relaciona a taxa anual de conversão da vegetação nativa para agricultura à interação entre o preço internacional defasado e a capacidade produtiva potencial do município:
\begin{equation}
Y_{it}^{\text{agr}} = \beta_{c}\,\big(\ln P_{c,t-1} \times Z_{i,c}\big) + \gamma_{c}\,\big(Z_{i,c} \times \tau_t\big) + \delta\,R_{it} + \alpha_i + \lambda_{b(i),t} + \varepsilon_{it},
\label{eq:modelo-principal}
\end{equation}
em que $i$ indexa o município e $t$ o ano da transição observada entre $t-1$ e $t$, de 2001 a 2024. A variável dependente $Y_{it}^{\text{agr}}$ mede a área de vegetação nativa convertida para agricultura no ano $t$ como porcentagem do estoque municipal de vegetação nativa disponível no início do período anual, expressa em pontos percentuais ($100 \times \text{taxa}$). O termo $\ln P_{c,t-1}$ corresponde ao logaritmo natural do preço real internacional da cultura $c$ no ano $t-1$, extraído da \textit{Pink Sheet} e deflacionado para dólares de 2024 pelo CPI-U \cite{WorldBank2026,BLSCPI2026}, centrado na média dos 24 anos da série. A defasagem situa a cotação no ano anterior à transição registrada, sem pressupor que essa ordenação anual elimine expectativas ou outros fatores comuns às decisões e aos preços.

A variável $Z_{i,c}$ representa o rendimento potencial municipal da cultura $c$ no GAEZ v4 sob regime de sequeiro e alto nível de insumos \cite{FAOGAEZ2021}, padronizado com média zero e desvio-padrão unitário entre os 141 municípios do estado. O coeficiente de interesse para H1 é $\beta_c$: ele compara a associação entre preço real defasado e conversão agrícola em municípios que diferem em um desvio-padrão de capacidade. Um valor positivo indica associação mais positiva onde o potencial é maior, sem medir o efeito absoluto da cotação em cada município.

Essa comparação considera diferenças permanentes entre municípios e mudanças anuais compartilhadas. O efeito fixo municipal $\alpha_i$ absorve características invariantes no tempo, incluindo o nível isolado de $Z_{i,c}$, a área territorial e a localização geográfica. O efeito fixo de bioma predominante por ano $\lambda_{b(i),t}$ absorve mudanças anuais comuns aos municípios de cada grupo (Amazônia, Cerrado e Pantanal), assim como o nível isolado de $\ln P_{c,t-1}$ e variações cambiais ou macroeconômicas comuns no ano $t$. A identificação estatística de $\beta_c$ provém, portanto, da interação entre preços e potencial, condicionada a essa estrutura.

Municípios com capacidades distintas também podem seguir trajetórias diferentes de conversão, mesmo sem responder às oscilações anuais das cotações. O termo $Z_{i,c} \times \tau_t$ representa essa possibilidade por uma tendência linear em décadas, $\tau_t = (t - 2001)/10$. A variável $R_{it}$ registra a anomalia anual de precipitação em centenas de milímetros em relação à média municipal de 1981--2000, considerando a relação das condições climáticas com o calendário agrícola e a detecção por satélite \cite{ChomitzThomas2003,AssuncaoEtAl2023}. O termo $\varepsilon_{it}$ é o erro idiossincrático, e os erros-padrão da referência são agrupados por município.

Os modelos de referência apresentam separadamente as exposições de soja e milho, cujos preços e potenciais são correlacionados. As séries anuais de preços reais apresentam correlação de $0{,}948$ nos 24 anos do painel, e os rendimentos potenciais padronizados, correlação transversal de $0{,}591$ entre os 141 municípios. Em uma comparação residual que remove efeitos de município e ano, a correlação entre $\ln P_{s,t-1} \times Z_{i,s}$ e $\ln P_{m,t-1} \times Z_{i,m}$ permanece elevada ($0{,}565$). Essa proximidade dificulta separar as associações específicas das duas culturas. As regressões individuais facilitam a leitura de cada exposição, mas não isolam um incentivo independente do outro; por isso, a estimação conjunta e o teste de Wald avaliam também a significância simultânea do bloco de grãos. A heterogeneidade territorial é examinada por interações contínuas com as frações municipais de área em cada bioma e por subamostras de predominância amazônica e cerradeira.

A forma de representar a capacidade merece atenção porque as condições biofísicas se traduzem em possibilidades produtivas sob tecnologias específicas \cite{LubowskiEtAl2008,BermanEtAl2023}. A referência de alto insumo é comparada com o rendimento potencial do GAEZ em sequeiro sob baixo insumo, o índice composto de adequação edafoclimática do GAEZ e o rendimento comercial médio da PAM entre 1991 e 2000, este último restrito aos municípios com ao menos sete anos de registro prévio \cite{IBGEPAM2026}. Na agricultura de larga escala em Mato Grosso, correção de solos, sementes melhoradas e insumos ampliaram as possibilidades de cultivo \cite{Spera2017}. Esse antecedente fundamenta a comparação entre cenários, cujo propósito é avaliar a sensibilidade da associação à medida de capacidade, e não identificar a tecnologia efetivamente adotada nas áreas convertidas.

\subsection{Potencial biofísico, ocupação pré-determinada e margem de intensificação}
\label{estrategia-margens}

O potencial do GAEZ e a ocupação agrícola anterior caracterizam dimensões distintas das alternativas territoriais. A primeira medida descreve condições produtivas padronizadas; a segunda informa quanto do município já era utilizado pela agricultura antes do painel. Onde há áreas ocupadas e oportunidades de ganho produtivo, a resposta aos preços pode envolver reorganização de usos e cultivo em duas safras, como discutem \citeonline{Angelsen2010}, \citeonline{MacedoEtAl2012} e \citeonline{Braganca2018}. Na base municipal, o potencial da soja no GAEZ e a participação agrícola média de 1985--2000 apresentam correlação de $-0{,}025$, o que reforça a pertinência de examiná-los separadamente.

Para examinar como a associação com os preços varia também conforme esse histórico, a especificação de referência incorpora medidas de ocupação agrícola anteriores ao período de análise:
\begin{equation}
\begin{aligned}
Y_{it}^{\text{agr}} &= \beta_{c}\,\big(\ln P_{c,t-1} \times Z_{i,c}\big) + \theta_{c}\,\big(\ln P_{c,t-1} \times H_{i}\big) \\
&\quad + \gamma_{c}\,\big(Z_{i,c} \times \tau_t\big) + \psi_{c}\,\big(H_{i} \times \tau_t\big) + \delta\,R_{it} + \alpha_i + \lambda_{b(i),t} + \varepsilon_{it},
\end{aligned}
\label{eq:modelo-margens}
\end{equation}
em que $H_i$ representa, alternativamente, a participação média da área agrícola no território municipal entre 1985 e 2000 ou a participação da soja em 2000, ambas provenientes do MapBiomas e padronizadas entre os 141 municípios. A primeira reúne as classes agrícolas e não inclui pastagens; a segunda corresponde à classe de soja. Enquanto $\beta_c$ descreve a associação diferencial ligada ao potencial, $\theta_c$ compara a associação entre preços e conversão conforme a ocupação agrícola prévia. Essa medida não representa toda a área anteriormente aberta, distinção necessária para interpretar a reorganização produtiva.

A expansão do milho oferece uma comparação complementar sobre as possibilidades de aproveitar áreas já utilizadas. Para examinar sua relação com a ocupação agrícola anterior, estima-se um modelo auxiliar cuja variável dependente é a área plantada de milho da PAM como porcentagem da área total municipal ($A_{it}^{\text{milho}}$):
\begin{equation}
\begin{aligned}
A_{it}^{\text{milho}} &= \beta_m^{A}\,\big(\ln P_{m,t-1} \times Z_{i,m}\big)
 + \phi\,\big(\ln P_{m,t-1} \times H_i\big) \\
&\quad + \gamma_m^{A}\,\big(Z_{i,m} \times \tau_t\big)
 + \kappa\,\big(H_i \times \tau_t\big) \\
&\quad + \delta\,R_{it} + \alpha_i + \lambda_{b(i),t} + u_{it}.
\end{aligned}
\label{eq:modelo-intensificacao}
\end{equation}
Nessa equação, $H_i$ corresponde à participação agrícola média de 1985--2000, e $\beta_m^{A}$ e $\gamma_m^{A}$ distinguem os coeficientes da exposição e da tendência de potencial no modelo de área. O parâmetro $\phi$ compara a associação com preços conforme a ocupação prévia, enquanto $\kappa$ mede o diferencial de trajetória por década associado a um desvio-padrão de $H_i$. A expansão das duas safras em Mato Grosso, documentada por \citeonline{VanWeyEtAl2013}, \citeonline{Spera2017} e \citeonline{SimoesEtAl2020}, dá fundamento econômico a essa comparação. A área municipal de milho, porém, não distingue safras nem localiza a expansão em parcelas anteriormente abertas. Sua leitura conjunta com $\theta_c$ oferece evidência indireta sobre o aproveitamento de áreas utilizadas, sem observar integralmente esse mecanismo.

\subsection{Custos logísticos de acesso a mercados e condicionamento institucional}
\label{estrategia-logistica-instituicoes}

A passagem do preço internacional à rentabilidade do produtor depende dos custos de produção e de escoamento \cite{Pfaff1999,ChomitzThomas2003,LubowskiEtAl2008}. Nos indicadores de fertilizantes e nas planilhas estaduais da CONAB utilizados no painel, mais de 95\% da variância concentra-se no tempo, acompanhando a elevação das cotações agrícolas a partir de meados dos anos 2000 \cite{CONAB2026}. A interação desses custos com a capacidade potencial pode, portanto, compartilhar variação com a exposição aos preços. Sua inclusão é tratada como sensibilidade exploratória, reconhecendo a dificuldade de separar movimentos conjuntos das séries, sem pressupor irrelevância econômica dos custos.

O acesso aos mercados acrescenta diferenças espaciais a essa comparação, pois a infraestrutura logística varia entre municípios e se alterou ao longo do período \cite{VanWeyEtAl2013,LathuilliereEtAl2014}. A partir dos centroides dos 141 municípios na malha do IBGE (SIRGAS 2000, EPSG:4674), três medidas distinguem a proximidade persistente aos portos do eixo Sul das mudanças associadas à disponibilidade de terminais. A primeira, $P_i^{\text{sul}}$, é a menor distância geodésica a Santos ou Paranaguá, padronizada e com sinal invertido, de modo que valores maiores representam maior proximidade.

A segunda medida, $D_{it}^{\text{min}}$, registra a distância geodésica mínima aos terminais considerados ativos em cada ano, em centenas de quilômetros, incorporando Miritituba (Itaituba), no Arco Norte a partir de 2014. A terceira combina a proximidade padronizada a Rondonópolis, $P_i^{\text{rond}}$, com o período posterior à inauguração do terminal ferroviário: $F_{it}^{2013} = P_i^{\text{rond}} \times \mathbb{I}[t \ge 2013]$. As variáveis representam acesso potencial, não frete ou trajeto realizado. Sua inclusão como controles ou em interação com $\ln P_{c,t-1} \times Z_{i,c}$ permite avaliar a estabilidade de $\beta_c$ e diferenças da associação conforme o acesso, sem tomar a inauguração como choque exógeno identificado.

As restrições administrativas podem modificar os custos esperados de abrir áreas e, por essa via, condicionar os incentivos \cite{North1991,Ostrom2009,Shimshack2014}. Para aproximar essa dimensão, o artigo compara registros de autuação e de embargo. Os autos classificados como Flora têm cobertura restrita aos pares com classificação temática completa ($N = 1.135$ na amostra efetiva de estimação entre 2002 e 2024), reduzindo a amostra em cerca de 60\% e alterando sua composição territorial \cite{IBAMA2026}. Essa limitação, descrita anteriormente, motiva o uso complementar dos termos de embargo lavrados pelo IBAMA entre 2000 e 2024.

Dos 12.284 termos únicos não cancelados registrados em Mato Grosso, 11.513 (93,7\%) possuem coordenadas válidas no filtro espacial. Os termos localizados por interseção ponto-no-polígono na malha do IBGE compõem o painel municipal. Na grade município--ano, a ausência de termo espacialmente atribuído recebe zero, e as áreas ausentes ou negativas foram tratadas como zero no processamento vigente. Esse procedimento permite manter a amostra de conversão, mas ausência de registro não certifica ausência de infração ou embargo.

Com essa medida, estimam-se duas especificações complementares na amostra observada completa ($N = 2.767$):
\begin{equation}
\begin{aligned}
Y_{it}^{\text{agr}} &= \beta_c\,\big(\ln P_{c,t-1} \times Z_{i,c}\big)
 + \rho\,\ln\big(1 + E_{i,t-1}^{\text{ha}}\big) \\
&\quad + \omega_c\,\big[(\ln P_{c,t-1} \times Z_{i,c})
 \times \ln(1 + E_{i,t-1}^{\text{ha}})\big] \\
&\quad + \gamma_c\,\big(Z_{i,c} \times \tau_t\big)
 + \delta\,R_{it} + \alpha_i + \lambda_{b(i),t} + \varepsilon_{it},
\end{aligned}
\label{eq:modelo-embargo-area}
\end{equation}
\begin{equation}
\begin{aligned}
Y_{it}^{\text{agr}} &= \beta_c\,\big(\ln P_{c,t-1} \times Z_{i,c}\big)
 + \eta_1\,\ln\big(1 + N_{i,t-1}^{\text{emb}}\big) \\
&\quad + \eta_2\,\big[\ln P_{c,t-1} \times Z_{i,c}
 \times \ln(1 + N_{i,t-1}^{\text{emb}})\big] \\
&\quad + \eta_3\,\big[\ln P_{c,t-1} \times \ln(1 + N_{i,t-1}^{\text{emb}})\big] \\
&\quad + \eta_4\,\big[Z_{i,c} \times \ln(1 + N_{i,t-1}^{\text{emb}})\big] \\
&\quad + \gamma_c\,\big(Z_{i,c} \times \tau_t\big)
 + \delta\,R_{it} + \alpha_i + \lambda_{b(i),t} + \varepsilon_{it},
\end{aligned}
\label{eq:modelo-embargo-mod}
\end{equation}
em que $E_{i,t-1}^{\text{ha}}$ é a soma das áreas registradas nos termos, em hectares, no município $i$ no ano $t-1$, sem representar uma união territorial livre de sobreposições, e $N_{i,t-1}^{\text{emb}}$ é a contagem de termos de embargo geocodificados lavrados no ano anterior. Na Equação~\eqref{eq:modelo-embargo-area}, $\rho$ descreve a associação com a área embargada quando a exposição preço--potencial é zero, enquanto $\omega_c$ permite que essa associação varie com a exposição. Na Equação~\eqref{eq:modelo-embargo-mod}, quando a contagem registrada é zero, $\beta_c$ representa a diferença na associação entre preços e conversão correspondente a um desvio-padrão adicional de capacidade. Esse gradiente preço--potencial varia por unidade adicional de $\ln(1+N^{\text{emb}})$ conforme $\eta_2$, condicionadas as interações inferiores representadas por $\eta_3$ e $\eta_4$.

Essas comparações aproximam o argumento de condicionamento institucional discutido por \citeonline{Sousa2016}, \citeonline{HardingEtAl2021} e \citeonline{AssuncaoEtAl2023}. Como os registros também podem responder às infrações, sua associação com a conversão não identifica um efeito dissuasório. As estimativas com autos de Flora são apresentadas paralelamente para comparar o comportamento dos indicadores sob a restrição de cobertura, sem atribuir toda diferença de resultado a essa restrição.

\subsection{Destinos comparativos, regimes temporais, inferência e limites de identificação}
\label{estrategia-inferencia}

A comparação com pastagem permite examinar se a associação diferencial de preços e potencial agrícola se restringe à abertura imediata para agricultura. Para isso, a Equação~\eqref{eq:modelo-principal} é reestimada com a taxa anual de conversão da vegetação nativa para pastagem ($Y_{it}^{\text{past}}$, em pontos percentuais da vegetação inicial), preservando a distinção entre destinos \cite{JungEtAl2022}. O GAEZ de soja e milho não mede produtividade pecuária; portanto, o resultado compara trajetórias de conversão conforme oportunidades agrícolas. Os padrões de reorganização territorial discutidos por \citeonline{MacedoEtAl2012} e \citeonline{SimoesEtAl2020} oferecem interpretações possíveis para essa comparação, sem que o coeficiente, isoladamente, identifique deslocamento entre usos.

A estabilidade temporal é relevante porque o painel reúne períodos distintos de preços e de governança ambiental. Entre 2004 e 2006, a implementação do Plano de Ação para Prevenção e Controle do Desmatamento na Amazônia Legal (PPCDAm) e da Moratória da Soja coincidiu com mudanças no patamar das cotações \cite{AssuncaoEtAl2015,HeilmayrEtAl2020}. Para avaliar quanto $\beta_c$ depende do contraste entre 2001--2005 e 2006--2024, o primeiro exercício restringe a referência à subjanela 2006--2024. O segundo substitui a tendência $Z_{i,c} \times \tau_t$ por degraus $Z_{i,c} \times \mathbb{I}[t \ge 2005]$ e $Z_{i,c} \times \mathbb{I}[t \ge 2006]$, examinando outra forma de representar diferenças entre períodos.

O terceiro exercício acrescenta à Equação~\eqref{eq:modelo-principal} a interação com o preço futuro, $\ln P_{c,t+1} \times Z_{i,c}$, como placebo temporal. Uma associação significativa com essa medida enfraquece a leitura do coeficiente defasado como resposta específica às flutuações anuais. O diagnóstico exige considerar persistência das séries, expectativas e outros fatores comuns, mas não distingue sozinho essas explicações nem identifica os efeitos das políticas que coincidem com a mudança de período.

A avaliação da precisão segue os critérios registrados para estes exercícios. Como há apenas 24 preços anuais e os municípios podem estar sujeitos a dependência espacial \cite{Conley1999}, o agrupamento por município é comparado com duplo agrupamento por município e ano e com agrupamento pelas 22 microrregiões geográficas do IBGE. Também se examinam efeitos fixos de microrregião por ano. Uma estimação em duas etapas utiliza as inclinações anuais nos 24 pontos temporais, com erros robustos à heterocedasticidade e à autocorrelação (Newey-West com duas defasagens), enquanto 24 reestimações excluem um ano de cada vez (\textit{leave-one-year-out}). Esses exercícios distinguem a persistência do sinal da precisão atribuída às diferentes fontes de variação.

A sensibilidade funcional e amostral é avaliada por winsorização de $Y_{it}^{\text{agr}}$ no percentil 99, transformação arco-seno hiperbólico $\text{asinh}(10\,Y_{it}^{\text{agr}})$ e exclusão dos quatro municípios de predominância pantaneira. A grade completa de 3.384 observações com zeros imputados é utilizada apenas como comparação, para documentar como a inclusão de valores sem conversão observada modifica as estimativas e sua precisão.

O desenho guarda relação com exposições interagidas da literatura de instrumentos do tipo Bartik e com preços ponderados pela aptidão \cite{GoldsmithPinkhamEtAl2020,BermanEtAl2023}, mas essa semelhança não estabelece as condições de identificação causal. O potencial do GAEZ é fixo no painel, incorpora informações territoriais de 2010 e pode se correlacionar com atributos geográficos e de infraestrutura que seguem dinâmicas próprias. A produção mato-grossense participa da oferta mundial de soja e milho, e a atuação fiscalizatória pode responder à pressão de desmatamento. Assim, as Equações~\eqref{eq:modelo-principal} a~\eqref{eq:modelo-embargo-mod} organizam associações condicionais e permitem avaliar sua estabilidade diante de medidas, controles e recortes distintos. A discussão dos resultados deve preservar esse alcance ao relacionar as estimativas aos mecanismos econômicos, sem tratá-las como efeitos causais isolados ou parâmetros de bem-estar.
